\section{Combustion Stability Analysis}

\subsection{Introduction}
This chapter details the physics-based models for combustion and feed system stability. The analysis covers low-frequency chugging, high-frequency acoustic modes, and feed system coupling (POGO/surge). The system uses both lumped-parameter and spatially-distributed wave propagation models to assess stability margins.

\subsection{Key Files Referenced}
The physics models documented in this chapter are implemented in:
\begin{itemize}
    \item \texttt{engine/pipeline/stability/analysis.py} - Main stability analysis routines
    \item \texttt{engine/pipeline/stability/spatial.py} - Spatially-distributed stability models
    \item \texttt{engine/pipeline/stability/enhanced.py} - Enhanced physics-based stability models
    \item \texttt{engine/pipeline/stability/physics.py} - Core stability physics relations
\end{itemize}

\subsection{Nomenclature}
\begin{description}
    \item[$f$] Frequency [\si{Hz}]
    \item[$\tau_{res}$] Residence time [\si{s}]
    \item[$L^*$] Characteristic length [\si{m}]
    \item[$c^*$] Characteristic velocity [\si{m/s}]
    \item[$a$] Speed of sound [\si{m/s}]
    \item[$V_c$] Chamber volume [\si{m^3}]
    \item[$A_t$] Throat area [\si{m^2}]
    \item[$Z$] Acoustic impedance [\si{Pa.s/m^3}]
    \item[$\alpha$] Wave growth rate [\si{1/s}]
    \item[$K$] Bulk modulus [\si{Pa}]
    \item[$\Delta P_{feed}$] Feed system pressure drop [\si{Pa}]
\end{description}

\subsection{Low-Frequency Stability (Chugging)}

\subsubsection{Residence Time Model}
The fundamental chugging frequency is estimated from the chamber residence time:
\begin{equation}
\tau_{res} = \frac{L^*}{c^*} = \frac{V_c}{A_t c^*}
\end{equation}
\begin{equation}
f_{res} = \frac{1}{2\pi \tau_{res}}
\end{equation}

\subsubsection{Helmholtz Resonance Model}
For chambers acting as a volume-compliance element, the Helmholtz frequency is:
\begin{equation}
f_H = \frac{a}{2\pi} \sqrt{\frac{A_t}{V_c L_{eff}}}
\end{equation}
where $L_{eff} \approx 0.5 d_t$ is the effective neck length.

\subsubsection{Stability Index}
A heuristic stability index is calculated based on chamber pressure, $L^*$, and frequency:
\begin{equation}
I_{stab} = \text{min}\left(1, \left(\frac{P_c}{P_{ref}}\right)^{0.3}\right) \cdot \text{min}\left(1, \left(\frac{L^*}{L^*_{ref}}\right)^{0.2}\right) \cdot f_{factor}
\end{equation}

\subsection{Acoustic Modes}

\subsubsection{Longitudinal Modes}
Longitudinal frequencies are modeled using an open-closed tube approximation:
\begin{equation}
f_{nL} = \frac{(2n-1) a}{4 L_c}
\end{equation}

\subsubsection{Transverse Modes}
Transverse (tangential/radial) modes for a cylindrical chamber:
\begin{equation}
f_{nT} = \frac{\alpha_n a}{\pi D_c}
\end{equation}
where $\alpha_n$ are the roots of the Bessel function derivative $J'_m(\alpha) = 0$ (e.g., 1.841, 3.054).

\subsection{Feed System Stability}

\subsubsection{POGO and Surge}
Feed line resonance frequencies:
\begin{equation}
f_{pogo} = \frac{a_{prop}}{4 L_{line}} \quad \text{(Closed-Open)}
\end{equation}
\begin{equation}
f_{surge} = \frac{a_{prop}}{2 L_{line}} \quad \text{(Closed-Closed)}
\end{equation}
where $a_{prop} = \sqrt{K/\rho}$.

\subsubsection{Water Hammer}
The theoretical maximum pressure spike for an instantaneous valve closure:
\begin{equation}
\Delta P_{wh} = \rho a_{prop} \Delta v
\end{equation}
The water hammer margin is defined as:
\begin{equation}
M_{wh} = \frac{\Delta P_{feed}}{\Delta P_{wh}}
\end{equation}

\subsection{Spatial Stability Analysis}
The enhanced model solves the wave equation spatially to determine local growth rates $\alpha$:
\begin{equation}
\alpha = \frac{\dot{E}_{in} - \dot{E}_{diss}}{2 E_{stored}}
\end{equation}
where:
\begin{itemize}
    \item $\dot{E}_{in} \propto P_c \dot{m} / \rho$ (Combustion energy input)
    \item $\dot{E}_{diss} \propto \rho u^2 / L_c$ (Viscous/acoustic dissipation)
    \item $E_{stored} = \frac{1}{2} \rho a^2$ (Acoustic energy density)
\end{itemize}

\subsection{Pintle-Specific Effects}
Pintle geometry introduces localized effects on stability:
\begin{itemize}
    \item \textbf{Impingement}: Fuel spray hitting the wall creates pressure fluctuation sources.
    \item \textbf{Recirculation}: Zones near the pintle tip provide either damping (via turbulence) or amplification.
    \item \textbf{Ablation}: Uneven wall recession creates impedance mismatches $Z = \rho a / A(x)$, reflecting waves.
\end{itemize}
